\begin{pastexam}[label={px:ntu-110-comm-dsa-43}]{臺大電信所 110 學年度（B）第 43 題}
（複選；每題 2.5 分，每答錯一個選項倒扣 0.5 分至該題 0 分止）

Consider the red-black tree in Figure 2, which of the following statements are true?\\
(Assume bottom-up insertion/deletion)
\begin{enumerate}[label=\Alph*.]
\item In Figure 2, after inserting key 14, then there are 6 red nodes.
\item In Figure 2, after inserting key 54, then key 49 and key 54 are siblings.
\item In Figure 2, after deleting key 22, then key 25 and key 73 are siblings.
\item In Figure 2, after deleting key 66, then key 35 is red.
\item In Figure 2, after deleting key 83, then key 89 is red.
\end{enumerate}
\begin{center}
\begin{tikzpicture}[b/.style={circle,draw,thick,inner sep=1pt,minimum size=6mm,font=\small\sffamily},
  r/.style={b,draw=red,dotted,very thick},>={Stealth[length=2mm]},
  eb/.style={thick,->},er/.style={very thick,red,dotted,->}]
\node[b] (58) at (0,0) {58};
\node[r] (35) at (-2.4,-1.1) {35}; \node[r] (73) at (2.4,-1.1) {73};
\node[b] (22) at (-3.45,-2.5) {22}; \node[b] (49) at (-1.35,-2.5) {49};
\node[b] (66) at (1.3,-2.5) {66}; \node[b] (83) at (3.5,-2.5) {83};
\node[r] (25) at (-2.85,-3.7) {25}; \node[r] (50) at (-0.75,-3.7) {50}; \node[r] (89) at (4.25,-3.7) {89};
\foreach \a/\b in {58/35,58/73,22/25,49/50,83/89} \draw[er] (\a) -- (\b);
\foreach \a/\b in {35/22,35/49,73/66,73/83} \draw[eb] (\a) -- (\b);
\end{tikzpicture}
\\*[4pt] Figure 2: Red-black tree
\end{center}
\end{pastexam}

\begin{pastsoln}
答案是 A、B、D。

圖中紅色節點畫成紅色虛線圈，指向紅色節點的邊也畫成紅色虛線。 原樹：根 $58$（黑），子節點 $35$、$73$
（紅）；$35$ 的子節點 $22$、$49$（黑），$73$ 的子節點 $66$、$83$（黑）；$25$、$50$、$89$（紅）分別是 $22$、$49$、
$83$ 的右子。 每條根到 NIL 的路徑都經過 $2$ 個黑節點（不含 NIL），是合法的紅黑樹。 以下每個選項都從原樹
開始做。

\textbf{A 對}。 插入 $14$：$14 < 58$、$< 35$、$< 22$，放 $22$ 的左子，塗紅。 父節點 $22$ 是黑的，不用調整
（資結書第~10.3.4~節）。 紅節點變成 $35, 73, 25, 50, 89, 14$，共 $6$ 個。

\textbf{B 對}。 插入 $54$：$54 < 58$、$> 35$、$> 49$、$> 50$，放 $50$ 的右子，塗紅。 父節點 $50$ 紅、叔叔
（$49$ 的左子）是 NIL（黑），而且 $54$、$50$ 同在右邊，是「直線」的情形：對祖父 $49$ 左旋，$50$ 塗黑、$49$
塗紅（資結書虛擬碼~10.6）。 $50$ 取代 $49$ 成為 $35$ 的右子，左子 $49$、右子 $54$，兩者是兄弟。

\textbf{C 錯}。 刪除 $22$：$22$ 是黑節點、只有一個紅色子節點 $25$。 用 $25$ 取代 $22$ 並塗黑，黑高度不變
（資結書第~10.3.7~節）。 $25$ 成為 $35$ 的左子，它的兄弟是 $49$；$73$ 是 $35$ 的兄弟，在上一層。

\textbf{D 對}。 刪除 $66$：$66$ 是黑色葉子，拿掉後 $73$ 的左邊少一個黑節點（雙重黑）。 兄弟 $83$ 是黑的，
$83$ 的\textbf{遠端}子節點（右子）$89$ 是紅的，屬於最後一種情形：對父節點 $73$ 左旋，$83$ 繼承 $73$ 原本的
顏色（紅），$73$ 和 $89$ 塗黑，結束（資結書虛擬碼~10.8）。 結果 $58$ 的右子變成 $83$（紅），子節點 $73$、
$89$（黑）。 調整只發生在 $73$ 這棵子樹，$35$ 仍是紅的。

\textbf{E 錯}。 刪除 $83$：$83$ 黑、只有一個紅色子節點 $89$，用 $89$ 取代並塗成\textbf{黑}色。 $89$ 變黑。
\end{pastsoln}
